\section{The applicable agreement changes which rights survive}
\label{sec:apple-award-versions}

The timing question can be narrowed by reading the later forms already
incorporated into Apple's filings. The annual exhibit list identifies a CEO
performance form effective 25 September 2022, filed as Exhibit 10.2 with the
quarter ended 31 December 2022. It also includes general and CEO performance
forms effective 29 September 2024 as Exhibits 10.20 and 10.22. These are
different versions and participant categories, even though their notices use
similar headings and leave the individual grant particulars blank
\citep[annual exhibit list; Exhibits 10.20 and 10.22]{apple2024}
\citep[Exhibit 10.2]{apple2023q1awards}.

The 2024 general form does not remove the earlier timing distinction.
Section 8 still provides that qualifying death or disability before vesting
allows the performance-qualified award, multiplied by the severance fraction,
to vest on the Vesting Date. Section 7 requires share delivery as soon as
administratively practical thereafter, within two and one-half months and
subject to its conditions. Thus neither inspected general form supplies a
rule to deliver the shares immediately on severance. The proxy's description
of accelerated vesting remains a summary whose application to a particular
award needs further evidence \citep[Exhibit 10.20, Sections 7--8]{apple2024}
\citep[printed p.~53]{apple2025proxy}. Later filing alone is not evidence
that an earlier grant was amended.

The CEO forms change the surviving entitlement. In both the September 2022
and September 2024 versions, Section 8(b) allows unvested units to remain
eligible for the Vesting Date following death, disability, or qualifying
retirement, subject to the Committee's stated authority. These clauses do not
apply the general form's severance fraction. Retirement combines an age of
at least 60 and at least ten years of service with the stated service scope;
the retirement exception normally also requires termination on or after the
first anniversary of the Award Date. Earlier retirement does not receive
that exception unless the Committee determines otherwise. Continued
eligibility still depends on performance and does not make the eventual
number of shares known on departure
\citep[Exhibit 10.2, Sections 3--4 and 8]{apple2023q1awards}
\citep[Exhibit 10.22, Sections 3--4 and 8]{apple2024}.

To isolate the effect, reuse the hypothetical 1,000-unit target and a
performance multiplier of 1.5, with no modification, recoupment or other
condition changing the result. Suppose death or disability occurs at a point
giving a general-form severance fraction of $3/5$. Under the general form,
$1{,}000\times1.5\times3/5=900$ units remain eligible to vest. Under the
CEO form's stated continuation provision, all $1{,}000\times1.5=1{,}500$
units remain eligible. With a negative-return cap reducing the multiplier to
one, the corresponding quantities are 600 and 1,000. These comparisons
hold the performance result fixed to expose the contract difference; they
are not facts about any executive's award.

The difference carries into cash rights without creating an earlier payment.
With the same hypothetical \$250 dividend-equivalent credit established
before departure, the uncapped gross payments are \$225 under the
general-form fraction and \$375 under the CEO continuation provision. The
units and cash rights remain linked to the same applicable terms and
performance conditions. The CEO forms' Section 7 anchors delivery to the
Vesting Date, and Section 5(b) ties dividend equivalents to the award's
vesting and payment terms. Moving those amounts to the severance date would
change the specified cash and share path
\citep[Exhibit 10.2, Sections 5(b), 7--8]{apple2023q1awards}
\citep[Exhibits 10.20 and 10.22, Sections 5(b), 7--8]{apple2024}.

A second boundary is independent of the performance calculation. An employee
aged 60 with ten years of service has not necessarily satisfied the CEO
form's first-award-anniversary condition. Conversely, passing that anniversary
does not establish the age and service conditions. The record therefore needs
the applicable agreement, Award Date, termination date and reason, service
history and any relevant Committee decision. Unknown values remain unknown;
they cannot be converted into automatic eligibility or forfeiture.

The later forms retain unfilled performance schedules and the same kinds of
return, comparison-group and adjustment dependencies. They do not fill the
rank convention or establish the actual October 2023 award terms by
themselves \citep[Exhibits 10.20 and 10.22, Exhibit A]{apple2024}. The
appropriate result of this comparison is a narrower conditional model: the
general and CEO forms support different survival rules, both preserve later
delivery, and the proxy's timing description still needs an award-specific
explanation. A source-reference match establishes which form was retained;
applicability establishes which form can govern a financial path.
